\subsection{Recorded Reassociation Program}
\label{app:reassociation-record}
This case connects the mathematical validation in
Proposition~\ref{prop:reassociation-program} to an emitted attribute program.
The exact rule and generated JSON are retained in
\nolinkurl{evaluation/reassociation/}. It omits all eight RHS axis fields
and uses the batch-preservation and integer-bound guard in
Appendix~\ref{app:batched-reassociation-program}, with the generalized dot contract
of Listing~\ref{lst:dot-declaration}.

\paragraph{Generated expressions.}
The compiler prints the outer replacement as \code{auto#2} and the inner
replacement as \code{auto#3}. Table~\ref{tab:reassociation-temporaries}
relates its temporaries to Equation~\eqref{eq:batched-program-locals}.
\code{match_positions} is a presentation of paired overlap selection;
its left/right fields are $u_G,v_G$. For the queries below, inclusion
proved in Proposition~\ref{prop:reassociation-program} makes
\code{matched_positions(query, domain)} equal to
$\pos(\mathit{domain},\mathit{query})$.

\begin{table}[!htbp]
\caption{Correspondence with the recorded compiler output.}
\label{tab:reassociation-temporaries}
\small
\begin{tabularx}{\linewidth}{@{}lY@{}}\toprule
Printed temporary & Paper expression\\\midrule
\code{m0}, \code{m2} & $i_X$, $j_Y$\\
\code{p0}, \code{p1} & $(u_I,v_I)$, $(u_J,v_J)$\\
\code{m4} & $b_Z\Vert z_J$\\
\code{m5} & $f_Y$\\
\code{m6} & $b+|f_Y|+\operatorname{range}(|f_Z|)$\\
\code{m7} & $\pos(f_Z,z_I)$\\
\code{m8} & $\gath(b+\operatorname{range}(|f_Y|),\pos(f_Y,k_Y))=s_K$\\\bottomrule
\end{tabularx}
\end{table}

In particular, the emitted outer right-contraction list is
\code{concat(gather(m6, m7), m8)}. The entrywise identity
\begin{equation}
 \gath(c+\operatorname{range}(n),q)=c+q
 \qquad\text{for an in-range selector }q
 \label{eq:shifted-range-selection}
\end{equation}
reduces it to $s_I\Vert s_K$, exactly the expression in
Table~\ref{tab:reassociation-program}. Both sides have entry $c+q_j$;
the guard bounds these integers. The remaining six direct source/selection
fields and the inner right-contraction field $z_J$ are as displayed in
that table. No inferred RHS field is treated as an input to another
attribute program.

\paragraph{Recorded execution.}
Semantic checking returned \code{Proved} without an additional guard;
compilation and portable validation succeeded. Exact sources, commands,
logs, hashes, and compiler provenance are in
\nolinkurl{evaluation/reassociation/}. The proof in
Appendix~\ref{app:reassociation-proof} justifies the coordinate bijection
and ordered output; the recorded trace supplies implementation evidence
and is not an independently checked proof certificate. This case study
is separate from the corpus totals.
